% Topic 1.3.04 — Matrix Multiplication and Numerical Linear Solvers.

\section{Matrix Multiplication and Numerical Linear Solvers}
\label{sec:1.3.04-matrix-numerical-linear}

\begin{ticket}
Representation of matrices and vectors. Matrix multiplication algorithms and efficient implementations. Numerical methods for solving systems of linear equations.
\end{ticket}

\subsection*{Theory}

This ticket combines a data-structure question (how to lay out
matrices and vectors in memory so that loops are fast) with two
algorithmic ones (multiplying matrices and solving $A \vect{x} = \vect{b}$).
For deeper treatment see \cite[Ch.~28]{CLRS2013}.

\subsubsection*{Representation}

A vector $\vect{v} \in \R^n$ is stored as a contiguous array of $n$
floating-point values; element access is $\Theta(1)$.

A matrix $A \in \R^{m \times n}$ has two principal dense layouts:
\begin{itemize}
  \item \emph{row-major}: $A[i][j]$ at memory offset $i \cdot n + j$
        (used by C, NumPy by default);
  \item \emph{column-major}: $A[i][j]$ at offset $j \cdot m + i$
        (used by Fortran, BLAS, MATLAB).
\end{itemize}
The choice matters for cache behaviour: a tight loop iterating along
the contiguous axis enjoys spatial locality, while iteration across
the strided axis incurs one cache miss per element in the worst case.

For a sparse matrix with $\mathrm{nnz}(A)$ non-zero entries, common
formats are:
\begin{itemize}
  \item \emph{COO} (coordinate): three arrays of length $\mathrm{nnz}(A)$
        storing $(i, j, A_{ij})$ for each non-zero;
  \item \emph{CSR} (compressed sparse row): $A_{ij}$ values and column
        indices grouped by row, with a length-$(m+1)$ row-pointer array;
        provides $\Theta(\mathrm{nnz}(A_{i,*}))$ row access.
\end{itemize}

\subsubsection*{Dense matrix multiplication}

Given $A \in \R^{n \times n}$, $B \in \R^{n \times n}$, compute
$C = AB \in \R^{n \times n}$ with $C_{ij} = \sum_{k=1}^{n} A_{ik} B_{kj}$.

\begin{proposition}[Naive matrix multiplication]
\label{prop:naive-matmul}
The triple-loop algorithm
\[
  \text{for } i, j, k \in \{1, \dots, n\}\colon\quad C_{ij} \mathrel{+}= A_{ik}\, B_{kj}
\]
computes $C = AB$ in $\Theta(n^3)$ time and $\Theta(n^2)$ extra space
(the matrix~$C$ itself).
\end{proposition}
\proofof{naive-matmul}

\begin{theorem}[Strassen 1969]
\label{thm:strassen}
There is an algorithm that multiplies two $n \times n$ matrices using
$T(n) = 7\, T(n/2) + \Theta(n^2)$ arithmetic operations, hence
$T(n) = \Theta(n^{\log_2 7}) = \Theta(n^{2.807\ldots})$.
\end{theorem}
\proofof{strassen}

\begin{remark}[Block multiplication and cache efficiency]
Even without changing the asymptotic exponent, partitioning $A$ and
$B$ into $b \times b$ blocks (with $b$ chosen so a block fits in
cache) gives a constant-factor speed-up over naive triple loops on
modern hardware: each loaded block is reused $\Theta(n/b)$ times
before eviction. With~$b = \Theta(\sqrt{M})$ where $M$ is the cache
size, the number of cache misses drops from $\Theta(n^3)$ to
$\Theta(n^3 / \sqrt{M})$.
\end{remark}

\subsubsection*{Direct solvers}

For $A \in \R^{n \times n}$ non-singular and $\vect{b} \in \R^n$,
\emph{Gaussian elimination} reduces $A$ to upper-triangular form by
forward elimination, then \emph{back-substitution} solves the
triangular system. Equivalently, the elimination factorises
$P A = L U$ where $L$ is unit lower-triangular, $U$ is upper-triangular,
and $P$ is a permutation matrix recording the row swaps used for
\emph{partial pivoting} (always pivot on the largest-magnitude entry
in the current column to limit numerical growth).

\begin{theorem}[Cost of LU with partial pivoting]
\label{thm:lu-cost}
Computing the factorisation $P A = L U$ for $A \in \R^{n \times n}$
takes $\tfrac{2}{3} n^3 + O(n^2)$ floating-point operations and
$O(n^2)$ extra memory. Given the factorisation, solving
$A \vect{x} = \vect{b}$ via $L \vect{y} = P \vect{b}$ followed by
$U \vect{x} = \vect{y}$ takes $2 n^2 + O(n)$ flops by triangular
solves.
\end{theorem}
\proofof{lu-cost}

For an SPD matrix, the analogous \emph{Cholesky factorisation}
$A = L L^\top$ exists, costs half as many flops, and avoids pivoting.

\subsubsection*{Iterative solvers}

For large sparse $A$, direct methods are dominated by the
$\Theta(n^3)$ cost of dense factorisation. Iterative methods produce
a sequence $\vect{x}^{(0)}, \vect{x}^{(1)}, \dots$ that converges to
the true solution.

A \emph{stationary iteration} is built from a splitting $A = M - N$
with $M$ easily invertible:
\[
  \vect{x}^{(k+1)} \;=\; M^{-1} \bigl(N \vect{x}^{(k)} + \vect{b}\bigr)
                      \;=\; T \vect{x}^{(k)} + \vect{c},
  \qquad T = M^{-1} N, \;\vect{c} = M^{-1} \vect{b}.
\]
Two classical choices, with $A = D + L_A + U_A$ for diagonal~$D$ and
strictly lower / upper triangular $L_A, U_A$:
\begin{itemize}
  \item \emph{Jacobi:} $M = D$, $N = -(L_A + U_A)$, so
        $T_J = -D^{-1}(L_A + U_A)$;
  \item \emph{Gauss--Seidel:} $M = D + L_A$, $N = -U_A$, so
        $T_{GS} = -(D + L_A)^{-1} U_A$.
\end{itemize}

\begin{theorem}[Convergence of stationary iterations]
\label{thm:stationary-convergence}
The iteration $\vect{x}^{(k+1)} = T \vect{x}^{(k)} + \vect{c}$
converges to the unique fixed point $\vect{x}^* = (I - T)^{-1} \vect{c}$
from every starting vector $\vect{x}^{(0)}$ if and only if the
\emph{spectral radius} $\rho(T) = \max\{|\lambda| : \lambda \text{ eigenvalue of } T\}$
is strictly less than~$1$. The asymptotic error reduction per step
is $\rho(T)$.
\end{theorem}
\proofof{stationary-convergence}

A common sufficient condition: if $A$ is strictly diagonally dominant
($|A_{ii}| > \sum_{j \neq i} |A_{ij}|$ for every~$i$), both Jacobi and
Gauss--Seidel converge.

\subsection*{Proofs}

\begin{proof}[\Cref{prop:naive-matmul}]
\proofanchor{naive-matmul}
The body $C_{ij} \mathrel{+}= A_{ik} B_{kj}$ executes once for each
$(i, j, k) \in \{1, \dots, n\}^3$, that is $n^3$ times; each
execution is $O(1)$. Initialising $C$ to zero costs $\Theta(n^2)$.
Total: $\Theta(n^3)$.
\end{proof}

\begin{proof}[\Cref{thm:strassen}]
\proofanchor{strassen}
Partition $A$, $B$, $C$ into $2 \times 2$ blocks of size $n/2 \times n/2$:
\[
  A = \begin{pmatrix} A_{11} & A_{12} \\ A_{21} & A_{22} \end{pmatrix},\quad
  B = \begin{pmatrix} B_{11} & B_{12} \\ B_{21} & B_{22} \end{pmatrix}.
\]
The naive blocked algorithm uses $8$ multiplications of $(n/2)$-sized
matrices, giving $T(n) = 8 T(n/2) + \Theta(n^2) = \Theta(n^3)$ — no
gain. Strassen gives explicit formulas for $C_{ij}$ using only
\emph{seven} block products $M_1, \dots, M_7$ of suitable additions/subtractions
of the $A_{ij}$ and $B_{ij}$, e.g.
\[
  M_1 = (A_{11} + A_{22})(B_{11} + B_{22}),\quad
  M_2 = (A_{21} + A_{22}) B_{11},\quad \dots
\]
The blocks $C_{ij}$ are then linear combinations of $M_1, \dots, M_7$.
The recurrence becomes $T(n) = 7 T(n/2) + \Theta(n^2)$, and by the
Master Theorem (\cref{thm:master}) with $a = 7, b = 2, \alpha = \log_2 7 \approx 2.807$
and $f(n) = \Theta(n^2) = O(n^{\alpha - \varepsilon})$ for any
$\varepsilon < \alpha - 2$, case~(a) gives $T(n) = \Theta(n^{\log_2 7})$.
\end{proof}

\begin{proof}[\Cref{thm:lu-cost}]
\proofanchor{lu-cost}
Eliminating column~$k$ ($k = 1, \dots, n - 1$) requires $n - k$
multipliers (entries below the diagonal in column~$k$ of~$L$) and
updates the remaining $(n - k) \times (n - k)$ submatrix with one
multiply and one subtract per entry: $(n - k)^2$ flops.

Total flops for forward elimination:
\[
  \sum_{k=1}^{n-1} \bigl[(n - k) + 2 (n - k)^2\bigr]
  \;=\; \tfrac{n(n - 1)}{2} + 2 \cdot \tfrac{(n - 1) n (2n - 1)}{6}
  \;=\; \tfrac{2}{3} n^3 + O(n^2).
\]
Triangular solves: forward solve $L \vect{y} = P \vect{b}$ uses
$\sum_{i=1}^{n} (i - 1)$ multiplications and the same number of
subtractions, $= n^2 + O(n)$ flops; back substitution
$U \vect{x} = \vect{y}$ uses another $n^2 + O(n)$. Total $2 n^2 + O(n)$.
\end{proof}

\begin{proof}[\Cref{thm:stationary-convergence}]
\proofanchor{stationary-convergence}
Let $\vect{e}^{(k)} = \vect{x}^{(k)} - \vect{x}^*$. Subtracting the
fixed-point equation $\vect{x}^* = T \vect{x}^* + \vect{c}$ from the
iteration gives $\vect{e}^{(k+1)} = T \vect{e}^{(k)}$, hence
$\vect{e}^{(k)} = T^k \vect{e}^{(0)}$.

\emph{Sufficiency.} If $\rho(T) < 1$, then $\|T^k\| \to 0$ in any
matrix norm (Gelfand's formula:
$\rho(T) = \lim_k \|T^k\|^{1/k}$, so for any
$r \in (\rho(T), 1)$ eventually $\|T^k\| \leq r^k \to 0$). Hence
$\vect{e}^{(k)} \to 0$ for every initial $\vect{e}^{(0)}$.

\emph{Necessity.} If $\rho(T) \geq 1$, take an eigenpair
$T \vect{v} = \lambda \vect{v}$ with $|\lambda| \geq 1$ and start
with $\vect{e}^{(0)} = \vect{v}$: then
$\|\vect{e}^{(k)}\| = |\lambda|^k \|\vect{v}\|$, which does not tend
to zero.

The asymptotic rate $\rho(T)$ follows from $\|T^k\|^{1/k} \to \rho(T)$.
\end{proof}

\subsection*{Examples}

\begin{example}[Naive vs. Strassen at $n = 1024$]
\label{ex:naive-vs-strassen}
Naive multiplication does $1024^3 \approx 1.07 \times 10^9$ scalar
multiplications. Strassen's recurrence applied recursively down to
some small base case~$n_0$ does roughly
$1024^{\log_2 7} \approx 1024^{2.807} \approx 1.32 \times 10^8$
scalar multiplications — about $8 \times$ fewer.

The crossover at which Strassen's constant overhead pays off is
typically $n_0 \approx 50$--$100$ on real hardware (Strassen's seven
block products are accompanied by 18 block additions/subtractions,
versus naive's 8 products and 4 additions). Production BLAS libraries
fall back to a tuned naive triple loop for $n \leq n_0$.
\end{example}

\begin{example}[LU on a $3 \times 3$ matrix]
\label{ex:lu-3x3}
Let $A = \begin{pmatrix} 2 & 1 & 1 \\ 4 & 3 & 3 \\ 8 & 7 & 9 \end{pmatrix}$.
Without pivoting:
\begin{enumerate}[label=(\arabic*)]
  \item subtract $\tfrac{4}{2} = 2$ times row~1 from row~2; row~2
        becomes $(0, 1, 1)$;
  \item subtract $\tfrac{8}{2} = 4$ times row~1 from row~3; row~3
        becomes $(0, 3, 5)$;
  \item subtract $\tfrac{3}{1} = 3$ times row~2 (now $(0, 1, 1)$)
        from row~3; row~3 becomes $(0, 0, 2)$.
\end{enumerate}
The factors are
$L = \begin{pmatrix} 1 & 0 & 0 \\ 2 & 1 & 0 \\ 4 & 3 & 1 \end{pmatrix}$,
$U = \begin{pmatrix} 2 & 1 & 1 \\ 0 & 1 & 1 \\ 0 & 0 & 2 \end{pmatrix}$.
A direct check verifies $LU = A$. Solving $A \vect{x} = (4, 11, 29)\T$
by forward then back substitution yields
$\vect{y} = (4, 3, 4)\T$ and $\vect{x} = (1, 1, 2)\T$.
\end{example}

\begin{example}[Jacobi on a small SPD system]
\label{ex:jacobi}
Let $A = \begin{pmatrix} 4 & 1 \\ 1 & 3 \end{pmatrix}$, $\vect{b} = (5, 4)\T$, true solution $\vect{x}^* = (1, 1)\T$.
The Jacobi iteration matrix is
$T_J = -D^{-1}(L_A + U_A) = -\begin{pmatrix} 1/4 & 0 \\ 0 & 1/3 \end{pmatrix}\begin{pmatrix} 0 & 1 \\ 1 & 0 \end{pmatrix} = \begin{pmatrix} 0 & -1/4 \\ -1/3 & 0 \end{pmatrix}$.
Its eigenvalues solve $\lambda^2 = 1/12$, so $\rho(T_J) = 1/\sqrt{12} \approx 0.289$
and the iteration converges, halving the error roughly every
$\log_2(1/\rho) \approx 1.79$ steps. Starting from $\vect{x}^{(0)} = 0$,
the first iterates are $(5/4, 4/3) \approx (1.25, 1.33)$,
$(5/4 - 4/12, 4/3 - 5/12) \approx (0.917, 0.917)$, converging
visibly toward $(1, 1)$.
\end{example}

\begin{problem}[Cache-aware blocking]
Suppose the cache holds $M$ floating-point values and matrices $A$, $B$,
$C$ are stored in row-major order. Show that the naive triple-loop
matrix multiplication of two $n \times n$ matrices incurs
$\Theta(n^3)$ cache misses when $n^2 > M$ (no useful row reuse). Then
show that block multiplication with block size $b = \Theta(\sqrt{M})$
incurs only $\Theta(n^3 / \sqrt{M})$ misses, achieving the
Hong--Kung lower bound for matrix multiplication.
\end{problem}
